\chapter{Introduction}

\section{Overview and Background}

A team that pushes code several times a day also breaks its build several times a day. Most of the breaks are small: a missing \texttt{require}, a misspelled method name, a status code compared the wrong way round. Each one still costs a person the same routine. They see the red mark on the commit, open the CI log, scroll to the first error, find the file, change a line, push again and wait for the pipeline to turn green. The fix is often shorter than the time spent finding it.

This project automates that routine. When a GitHub Actions run fails with a code error, a workflow collects the failing job's log and the source files involved, asks a language model what is wrong, and opens a pull request with the change. A person still reviews and merges the pull request, and the system never pushes to the main branch.

The repository holds four cooperating parts, which Chapter 4 describes in detail.

\begin{itemize}
  \item \textbf{GitHub Actions workflow:} a workflow for a small Express demo application. When its tests fail with a code error, a notify job calls a webhook (source: \texttt{.github/\allowbreak{}workflows/\allowbreak{}demo-ci.yml}).
  \item \textbf{n8n workflow:} 26 nodes that receive the call, gather the CI log and source files through the GitHub API, build the prompt, call the model, apply the reply and open the pull request (source: \texttt{n8n/\allowbreak{}self\_\allowbreak{}healing\_\allowbreak{}pipeline.json}).
  \item \textbf{Model endpoint:} Qwen3.8-27B served by llama.cpp on a Kaggle notebook with two T4 GPUs and published through a Cloudflare quick tunnel (source: \texttt{kaggle/\allowbreak{}endpoint/\allowbreak{}deploy\_\allowbreak{}src.py}).
  \item \textbf{ci-heal-backend:} a second healer written as a FastAPI service. It classifies a failure first and then re-runs the job, opens a fix pull request or escalates to an issue. It talks to the same model endpoint (source: \texttt{ci-heal-backend/\allowbreak{}README.md}).
\end{itemize}

Around these sit double-click launcher scripts, an evaluation harness with its measured results, and a test rig that runs the workflow's real code inside a real n8n engine.

Three constraints shaped most of the later decisions.

\begin{itemize}
  \item \textbf{No paid model API:} the model runs on a free Kaggle notebook. That limits the work to a quantized model of roughly 27 billion parameters, a session that the notebook caps at 1.5 hours, and a network address that changes on every start.
  \item \textbf{The model is not trusted:} its reply is parsed, applied to the file by exact text matching, syntax-checked and tested by CI on the pull request before anyone looks at it. A reply that fails a step becomes a comment that gives the reason, not a pull request.
  \item \textbf{One healer per repository:} two healers reacting to the same failed run would open two competing pull requests, so the design assumes that exactly one of them is switched on.
\end{itemize}

\section{Problem Statement}

Fixing a small build break by hand is cheap in thought and expensive in attention. The developer has to leave what they were doing, read a log written for machines and make a change that is usually obvious once found. The aim here is to take over the finding and the first draft of the change, and to leave the decision to merge with a person.

Doing this safely is harder than calling a model with a log. Four problems had to be solved, and each one appears again in the later chapters.

\begin{itemize}
  \item \textbf{Not every red run is a code bug.} A busy port or a failed install has nothing to fix in the source. If the model were asked anyway, it would invent a patch. The demo tests therefore exit with 75 for environment problems, and the workflow maps that code to \texttt{infra} and sends nothing to the model (source: \texttt{demo-app/\allowbreak{}test.js}, \texttt{.github/\allowbreak{}workflows/\allowbreak{}demo-ci.yml}).
  \item \textbf{The model can be wrong.} Its reply may not match the file, may not parse, or may parse and still fail the tests. Every step after the reply is therefore a gate, and a fix that CI rejects can be retried a limited number of times.
  \item \textbf{Free hardware is slow and unstable.} A baseline run measured 12.29 tokens per second, so the first design, which asked for the whole fixed file, was slow. The tunnel in front of the notebook cuts a request after 120 seconds of silence. The address changes on every boot (source: \texttt{docs/\allowbreak{}results/\allowbreak{}DIFF\_\allowbreak{}FORMAT\_\allowbreak{}EVAL.md}, \texttt{docs/\allowbreak{}results/\allowbreak{}GOLIVE\_\allowbreak{}RESULTS.md}).
  \item \textbf{Claims need evidence.} It is easy to say that such a pipeline works. The goal was to measure how often the proposed fix is right, how long it takes from push to pull request, and what happens when it is wrong.
\end{itemize}

\section{Objectives of the System}

The system has eight objectives. Each has a short identifier so that Chapter 12 can trace it to the evidence that supports or limits it.

\begin{itemize}
  \item \textbf{O1:} Detect code failures and ignore the rest. A failed CI run with a code error must start the healer. A run that failed for an environment reason must start nothing and cause no model call.
  \item \textbf{O2:} Collect the evidence automatically. The healer must gather the failed job's log, the changed files and the source of the candidate files, and build one prompt from them without human input.
  \item \textbf{O3:} Obtain a small, checkable fix from a self-hosted model. The model must answer with SEARCH/REPLACE edits that can be applied by exact text matching, on hardware that costs nothing to use.
  \item \textbf{O4:} Deliver the fix as a pull request and never as a push to main. The fix must go to an \texttt{ai-fix-} branch and be reviewed and merged by a person.
  \item \textbf{O5:} Retry within a fixed budget. When CI fails on the fix pull request, the healer must try again with the model's reasoning switched on, up to a cap, and then stop and say so.
  \item \textbf{O6:} Offer a second, classifying healer. ci-heal-backend must classify a failure and then re-run, fix or escalate, with named thresholds and an attempt budget, and must write nothing to GitHub in its default mode.
  \item \textbf{O7:} Measure honestly. Fix quality, push-to-pull-request time and failure behaviour must be measured on real runs, with the decision rules written down before the data was collected where possible.
  \item \textbf{O8:} Keep it safe, testable and runnable by one person. Secrets must stay out of the repository, the logic must be testable without a GPU, and each component must start with a launcher script.
\end{itemize}

\section{Scope and Limitations}

\subsection{In scope}

The work covers one demo application, a small Express app with one to three source files, and failures that show up as failing tests after a push to \texttt{main} or on an \texttt{ai-fix-} pull request. It covers the full chain from the failed run to the open pull request, the model endpoint, the second healer, the evaluation harness and the test layers.

\subsection{Limitations}

The measured results are real but narrow. The main limits are listed here, and Chapter 14 of the project documentation holds the full list.

\begin{itemize}
  \item \textbf{One demo application:} all heal measurements use the same small Express app. Large code bases, other languages and build configuration are not covered.
  \item \textbf{Small case sets:} the quality studies use 11 to 17 cases with two or three repetitions each, so a difference of one case is within noise.
  \item \textbf{First-try quality is partial:} with thinking off, 12 of 18 hard single-file bugs were healed. With thinking on, 18 of 18 were, at about four times the time (source: \texttt{docs/\allowbreak{}results/\allowbreak{}SHIPPED\_\allowbreak{}QUALITY\_\allowbreak{}EVAL.md}).
  \item \textbf{Multi-file healing is below its own bar:} the first pass healed 9 of 18 against a bar of 12, and the retry loop reached 15 of 18 (source: \texttt{docs/\allowbreak{}results/\allowbreak{}MULTIFILE\_\allowbreak{}EVAL.md}). The report does not call multi-file healing finished.
  \item \textbf{A passing fix can still be wrong:} the gates check that the edit applies, parses and passes the repository's tests. In the recorded run, the fix removed the error and dropped a string conversion, and CI passed anyway.
  \item \textbf{The OpenAI fallback has never produced a fix:} its routing is tested, but a real fallback answer was never observed.
  \item \textbf{Failures that are not test failures are not healed:} a failed install, or a build step that fails after the tests pass, is ignored (source: \texttt{.github/\allowbreak{}workflows/\allowbreak{}demo-ci.yml}).
  \item \textbf{Operational limits:} the model address changes on every Kaggle boot, a session waits in the GPU queue without a bound, the n8n Cloud instance is a 14-day trial, seven workflow nodes have no error setting, and 17 of 21 backend audit findings are open (source: \texttt{docs/\allowbreak{}results/\allowbreak{}BACKEND\_\allowbreak{}AUDIT.md}).
\end{itemize}

\section{Organisation of the Report}

The report has 12 chapters and one appendix.

\begin{itemize}
  \item \textbf{Chapter 1, Introduction:} the problem, the objectives O1 to O8 and the scope.
  \item \textbf{Chapter 2, Literature Review:} the related work and background that the design builds on.
  \item \textbf{Chapter 3, System Requirements:} the roles, the functional and non-functional requirements and the use cases.
  \item \textbf{Chapter 4, System Architecture:} the components, the data flow of one heal, the trust boundaries and the main design decisions.
  \item \textbf{Chapter 5, GitHub Actions and the Notify Job:} the workflow file, the failure classification and the attempt counter.
  \item \textbf{Chapter 6, The n8n Healing Workflow:} the 26 nodes, stage by stage, and the retry behaviour.
  \item \textbf{Chapter 7, The Model Endpoint:} the Kaggle notebook, the server flags and the tunnel.
  \item \textbf{Chapter 8, The ci-heal-backend Service:} the classifier, the decision rules, the safety rules and the audit.
  \item \textbf{Chapter 9, Evaluation and Results:} the measured studies and the recorded live run.
  \item \textbf{Chapter 10, Security, Testing and Deployment:} credentials, the gates around the model, the test layers and how to run the system.
  \item \textbf{Chapter 11, Standards Adopted:} the design, coding and testing standards and where the repository enforces them.
  \item \textbf{Chapter 12, Conclusion:} the results traced back to O1 to O8, the open problems and future work.
  \item \textbf{Appendix A:} the repository map and directory structure.
\end{itemize}
